\documentclass[12pt,a4paper]{report}

\usepackage[margin=1in]{geometry}
\usepackage{graphicx}
\usepackage{booktabs}
\usepackage{longtable}
\usepackage{xcolor}
\usepackage{hyperref}
\usepackage{float}
\usepackage{array}
\usepackage{enumitem}

\hypersetup{
    colorlinks=true,
    linkcolor=blue,
    citecolor=blue,
    urlcolor=blue
}

\begin{document}

%========================
% TITLE PAGE
%========================

\begin{titlepage}
\centering

\vspace*{2cm}

{\Huge \textbf{ManuVision AI}}\\[0.5cm]

{\Large Manufacturing Vision AI Model Evolution Report}\\[1cm]

\textbf{Project Objective}

\\[0.3cm]

Design and continuously improve an AI-based manufacturing defect detection system through seven iterative model versions.

\vfill

\begin{tabular}{rl}
Author: & Your Name \\
Project Duration: & 2026 \\
Repository: & GitHub Link \\
Version: & Master Engineering Report
\end{tabular}

\vfill

\today

\end{titlepage}

\tableofcontents
\listoffigures
\listoftables

\chapter{Project Overview}
%\input{chapters/00_project_overview}
\section{Problem Statement}

Manufacturing industries require automatic visual inspection of components to identify defective products with high accuracy and low inference latency.

\section{Project Goal}

Develop a computer vision system capable of detecting surface defects and progressively improve its robustness across seven engineering iterations.

\section{Final Success Criteria}

\begin{itemize}
    \item High defect detection accuracy
    \item Low false negative rate
    \item Robust to illumination variation
    \item Robust to material variation
    \item Production-ready inference speed
\end{itemize}

\section{Development Philosophy}

Each version solves one or more identified limitations from the previous version. No modification is introduced without documenting:

\begin{enumerate}
    \item Existing limitation
    \item Root cause
    \item Hypothesis
    \item Engineering modification
    \item Experimental result
    \item Remaining limitations
\end{enumerate}

\chapter{Base Model}
%\input{chapters/01_base_model}

\chapter{Version 1}
%\input{chapters/02_v1}

\chapter{Version 2}
%\input{chapters/03_v2}

\chapter{Version 3}
%\input{chapters/04_v3}

\chapter{Version 4}
%\input{chapters/05_v4}

\chapter{Version 5}
%\input{chapters/06_v5}

\chapter{Version 6}
%\input{chapters/07_v6}

\chapter{Version 7}
%\input{chapters/08_v7}

\chapter{Final Comparative Analysis}

\chapter{Future Work}

\end{document}